% 第 6 题：理想变压器电路图（circuitikz）
% 拓扑（由解析反推，自洽）：原边 C--L1--原线圈--D；副边两并联支路：
%   X--L2--滑片 P 与 X--a--(变阻器 a-P 段) 及 (P-b 段)--b--Y、P--R0--Y；
% 滑片在中点时副边负载电阻最大。
\begin{tikzpicture}[font=\small, scale=0.9, every node/.style={transform shape=false}]
  % 变压器
  \draw (0,0) node[transformer core] (T) {};
  % 原边：C 端 -- L1 -- 原线圈上端；下端 -- D 端
  \draw (-2.2,0.5) to[lamp, l_=$L_{1}$] (T.A1);
  \draw (T.A2) -- (-2.2,-0.5);
  \node[circ, label=above:$C$] at (-2.2,0.5) {};
  \node[circ, label=below:$D$] at (-2.2,-0.5) {};
  % 副边：X = T.B1，Y = T.B2
  \coordinate (X) at (T.B1);
  \coordinate (Y) at (T.B2);
  % 支路 1：X -- L2 -- 滑片 P
  \draw (X) to[lamp, l=$L_{2}$] (2.6,0.55);
  \coordinate (P) at (2.6,0.55);
  % 变阻器（a 上端、b 下端），滑片位于 P
  \draw (X) -- (2.6,1.6);
  \draw (2.6,1.6) to[potentiometer, l=$R$] (2.6,-0.6);
  \coordinate (b) at (2.6,-0.6);
  \draw (b) -- (Y);
  % 支路 2：P -- R0 -- Y
  \draw (P) to[R, l_=$R_{0}$] (1.2,-0.5) -- (Y);
  % 触点标注
  \node[above left] at (2.6,1.6) {$a$};
  \node[below left] at (2.6,-0.6) {$b$};
  % 滑片连线到 P（potentiometer 滑片默认朝右，这里补一条短线到 P）
  \draw (2.6,0.55) -- (2.85,0.55);
\end{tikzpicture}
